\documentclass[10pt,twocolumn]{article}
\usepackage[margin=0.75in]{geometry}
\usepackage{booktabs}
\usepackage{amsmath}
\usepackage{graphicx}
\usepackage{siunitx}
\usepackage{microtype}
\usepackage{hyperref}
\usepackage{xcolor}
\usepackage{listings}
\lstset{basicstyle=\ttfamily\footnotesize,breaklines=true,columns=fullflexible}
\sisetup{group-separator={,},group-minimum-digits=4}
\hypersetup{colorlinks,linkcolor=blue!60!black,urlcolor=blue!60!black,citecolor=blue!60!black}

\title{RTL Simulation of Gemma~3 270M Inference on the Coral NPU:\\
Test Design, Numerics, Cycle Accounting and Simulator Throughput}
\author{Rachit Tibrewal \\ \small ergodex / coralnpu working tree, branch \texttt{feature/model-configurable-npu}}
\date{5 September 2026}

\begin{document}
\maketitle

\begin{abstract}
We report the first end-to-end, cycle-accurate RTL simulation of Gemma~3 270M
decoder inference on the Coral NPU (\texttt{RvvCoreMiniHighmemAxi}: RV32IMF +
Zve32f, VLEN\,=\,128, 1\,MB ITCM, 1\,MB DTCM). A six-tier cocotb test package
composes the repository's per-kernel RVV Gemma kernels into a full decoder
layer, an 18-layer forward pass and greedy generation, checked against a
bit-faithful NumPy replay and against the bf16 HuggingFace model. With real
weights, one decoder layer matches the strict reference to a relative error of
$2.6\times10^{-9}$ on the residual stream and the HF bf16 hidden state at
cosine 0.999995. One layer costs \num{2132836} cycles at 2.61 int8 MAC/cycle,
93\,\% of which is the int8 GEMV path, so a full token is $\approx$38\,M
cycles. The initial simulator throughput of 7.2\,k cycles/s was set entirely by
Python-in-the-loop bus modelling; moving external memory into the RTL behind a
DPI backdoor and removing per-cycle cocotb wake-ups raises it to
16.5\,k cycles/s single-threaded and 24.8\,k cycles/s with
Verilator threads, i.e.\ about 22 minutes per token instead
of 90. We give the resulting per-token latency
projections for the 50\,MHz FPGA and the 1\,GHz datasheet target and the
implications for the GR-288 accelerator work.
\end{abstract}

\section{Introduction}
The Coral NPU repository ships RVV kernels for every Gemma~3 building block
(int8/fp32/bf16 GEMV and matmul, FlashAttention, RMSNorm, tanh-GELU$\times$up,
residual add) and cocotb tests that check each in isolation against NumPy
golden values on tensors dumped from HuggingFace. What did not exist was a
test that runs the \emph{model}: weights for all 18 layers in external memory,
a KV cache that grows across tokens, RoPE at non-zero positions, and a
next-token prediction compared with the reference model. This report describes
that test package (\texttt{tests/cocotb/rvv/ml\_ops/gemma\_inference/}), the
numerical results obtained on an EC2 dev host, the cycle budget of one token on
the baseline core, and the simulator engineering needed to make multi-token
runs practical.

Gemma~3 270M (\texttt{config.json}, verified) has 18 decoder layers, hidden
size 640, FFN width 2048, 4 query heads and 1 KV head of dimension 256,
\texttt{query\_pre\_attn\_scalar}\,=\,256, a 262\,144-entry vocabulary tied to
the embedding table, sliding-window attention (512) on five of every six layers
and global attention on layers 5, 11 and 17, RoPE bases of $10^4$ (local) and
$10^6$ (global), Gemma-style RMSNorm $x\cdot\mathrm{rsqrt}(\overline{x^2}+\epsilon)\cdot(1+w)$
with $\epsilon=10^{-6}$, q/k head norms, and \texttt{gelu\_pytorch\_tanh}. No
logit softcapping is configured.

\section{Numerics under test}
Every projection ($W_q,W_k,W_v,W_o,W_{gate},W_{up},W_{down}$ and the lm\_head)
is executed as int8$\times$int8$\to$int32 on \texttt{rvv\_gemv\_int8}. Weights
are quantised per output channel (symmetric, round-to-nearest-even, clipped to
$\pm127$); activations are quantised per token at run time with the same rule.
The kernel's contract that $-128$ never appears on both operands is satisfied by
construction. Everything else is fp32: norms, RoPE, attention (the existing
unmasked FlashAttention kernel with a Taylor exponential), GELU (rational tanh),
the residual stream, the K/V caches and the logits. Prefill is executed as
sequential decode: with $Q_{\text{len}}=1$ over a cache holding positions
$0..p$, the unmasked kernel is exactly causal, which removes the need for a
masked prefill kernel at the cost of throughput.

Two references are used. The \emph{strict} reference (\texttt{gemma3\_ref.py})
replays the runner's dataflow in NumPy with the same int8 weights and
quantiser and exact $\exp$/$\tanh$; int8 GEMV results are bit-exact and every
fp32 stage is checked on a tap (post-RoPE $q$, attention output,
$\mathrm{gelu}\times\mathrm{up}$, residual, logits). The \emph{HF bf16}
reference (\texttt{dump\_gemma3\_model.py}) runs the model in its native
bfloat16 and stores hidden states after every layer, the last-position logits,
per-position top-$k$ ids and the greedy continuation. The strict reference
attributes a failure to one stage; the HF reference measures the quantisation
scheme itself, which is the gate the project's numerics rule requires.

\section{Test architecture}
\paragraph{DUT program.} \texttt{gemma3\_runner.cc} executes one command per
run for one token position: \texttt{CMD\_LAYERS} (a layer range on the
residual vector), \texttt{CMD\_LM\_HEAD} (final norm + lm\_head + argmax) or
\texttt{CMD\_FORWARD}. It never allocates weights: a 431-word descriptor in
DTCM holds a DDR address for every tensor (\texttt{gemma3\_model.h}), so the
ELF is 13.6\,KB of code and independent of the model image, the candidate
vocabulary or the sequence length. \texttt{mcycle} is read around every kernel
call and accumulated into twelve op buckets.

\paragraph{Memory image.} The test packs all weights into one 96.8\,MiB image
(\texttt{ModelImage}) and writes it directly into the testbench memory model:
at 16\,B per AXI beat, loading it over the bus would take longer than the
inference. The core fetches weights through its AXI master exactly as it would
from DDR. The 168\,MB embedding table stays host-side; the test supplies the
scaled embedding vector per token and scores the lm\_head on a chosen column
set ($\le 4096$ per run, 530 candidates in the default configuration, the whole
vocabulary in 64 chunks for tier~6).

\paragraph{Tiers.} T1 one layer at position 0 (every kernel once); T2 each
layer in isolation, fed HF's residual stream at its input (localises a failure,
covers both RoPE bases); T3 layer 0 over positions $0..N{-}1$ (cache growth,
RoPE at $p>0$, softmax over $N$ keys); T4 full forward over the prompt with
candidate logits and a top-1 check against HF; T5 greedy generation compared
token-by-token; T6 full-vocabulary logits.

\section{Results with real weights}
\subsection{Accuracy}
Table~\ref{tab:acc} lists the T1 taps for layer~0 at position~0 on the prompt
``The capital of France is''. The int8 GEMV stages are exact; the attention
output agrees to fp32 rounding because a single key makes the softmax trivial;
the only stage with visible error is the rational-tanh GELU
($4.8\times10^{-4}$ relative to the maximum). The residual stream after the
layer matches the strict reference to $2.6\times10^{-9}$ and the HF bf16
hidden state at cosine 0.999995. On a synthetic model with $\mathcal N(0,0.02)$
weights the GELU error was 20$\times$ larger ($1.0\times10^{-2}$), which
bounds the approximation on adversarial inputs.
With external memory in the RTL the results are bit-identical
(Table~\ref{tab:acc} was reproduced on both models).

The first multi-position runs taught two things. First, the strict reference
must replay the kernels' approximations, not exact functions: the
FlashAttention exponential is a range-reduced cubic Taylor polynomial
($8\times10^{-4}$ relative) and the GELU tanh is $y(y^2+27)/(9y^2+27)$ on
$\mathrm{clamp}(z,-3,3)$ ($2.4\times10^{-2}$ absolute). Because the FFN
intermediate is int8-quantised per tensor with Gemma's outliers, a
$10^{-3}$ difference in GELU flips int8 codes whose step is large, and the
residual diverged by $3.6\times10^{-2}$ at prefill position~2 although the
attention output matched to $3\times10^{-5}$. With the kernel arithmetic in
the reference, T3 passes all six prefill positions with every tap within
$\sim\!10^{-7}$ and HF bf16 cosines between 0.9972 and 0.9995. Second,
HuggingFace's last entry in \texttt{output\_hidden\_states} is taken after
the final RMSNorm, so the last-layer residual is compared after the final
norm; the raw comparison reads as cosine 0.15 and looks like a broken layer.
With both fixes, T2 passes: each of the 18 layers, fed HF's residual at its
input, matches the strict reference at cosine 1.000000 (relative error
$8\times10^{-8}$ to $1.2\times10^{-3}$) and HF bf16 at cosine $\geq 0.99999$,
on both RoPE bases. The first full 18-layer forward matches the strict
reference to $1.2\times10^{-7}$ on the residual and $6.6\times10^{-7}$ on the
candidate logits, so the DUT executes the int8 model exactly; the int8 model
itself drifts from bf16 HF to cosine 0.9875 after the final norm, the
accumulated cost of per-tensor activation quantisation. A complete token
(18 layers plus a 530-column lm\_head) measures \num{32011264} cycles at
3.14 MAC/cycle, 21 minutes of wall time on the two-thread model; the first
predicted token agrees with the strict reference. The top-1 comparison over
the six-token prompt was running when this report was written.

\begin{table}[t]
\centering\small
\caption{T1, layer 0, position 0, real weights. ``strict'' is the NumPy replay,
``HF'' the bf16 model.}
\label{tab:acc}
\begin{tabular}{lccc}
\toprule
Tap & cos (strict) & rel.\ max err & cos (HF) \\
\midrule
$q$ after norm+RoPE & 1.000000 & $1.1\times10^{-7}$ & -- \\
attention output & 1.000000 & $6.0\times10^{-8}$ & -- \\
gelu(gate)$\times$up & 0.999995 & $4.8\times10^{-4}$ & -- \\
residual out & 1.000000 & $2.6\times10^{-9}$ & 0.999995 \\
\bottomrule
\end{tabular}
\end{table}

\subsection{Cycle accounting}
One decoder layer takes \num{2132836} cycles for \num{5570560} int8 MACs,
2.61 MAC/cycle, with the Python AXI memory model (Table~\ref{tab:cyc}); with
the one-cycle RTL memory of Section~6 the same layer takes \num{1780644}
cycles (3.13 MAC/cycle), which shows that about 17\,\% of the original
layer time was external-memory latency rather than compute. The three GEMV buckets are 93\,\% of the
layer; the fp32 glue (norms, quantise, dequantise, RoPE, GELU, residual) is
4.2\,\% and single-key attention 0.5\,\%. Multiplying out, an 18-layer token is
32 to 38\,M cycles before the lm\_head depending on memory latency, and the lm\_head over the full
vocabulary (168\,M MACs) would add a further $\approx 64$\,M cycles at the same
rate: for a decode step the tied embedding matrix costs more than the
transformer body.

\begin{table}[t]
\centering\small
\caption{Per-op cycles, layer 0, real weights (int8 MAC/cycle in
parentheses).}
\label{tab:cyc}
\begin{tabular}{lrr}
\toprule
Op & Cycles & Share \\
\midrule
gate/up GEMV (2$\times$640$\times$2048) & \num{933844} & 43.8\,\% (2.81) \\
down GEMV (2048$\times$640) & \num{466262} & 21.9\,\% (2.81) \\
q/k/v GEMV (640$\times$1536) & \num{350286} & 16.4\,\% (2.81) \\
o GEMV (1024$\times$640) & \num{233302} & 10.9\,\% (2.81) \\
RMSNorm ($\times$7) & \num{31711} & 1.5\,\% \\
gelu$\times$up & \num{15317} & 0.7\,\% \\
dequantise & \num{14173} & 0.7\,\% \\
quantise & \num{12747} & 0.6\,\% \\
residual add ($\times$2) & \num{11605} & 0.5\,\% \\
attention (1 key) & \num{9777} & 0.5\,\% \\
RoPE & \num{3397} & 0.2\,\% \\
\midrule
total & \num{2132836} & 100\,\% (2.61) \\
\bottomrule
\end{tabular}
\end{table}

The GEMV rate of 2.8 MAC/cycle is far from the vector unit's peak (about 32
int8 ALU lanes per cycle per unit pair): with VLEN\,=\,128, one \texttt{e8m4}
group is 64 bytes, the reduction per output column serialises on
\texttt{vredsum}, and every weight byte crosses the single 16\,B/cycle LSU port
from external memory. Weight bandwidth alone bounds a layer at
$5.57\,\text{M}/16 = 348$\,k cycles, so the kernel is compute/issue bound, not
memory bound, at this size.

\subsection{Projected latency}
Taking the one-cycle-memory figure, a token is 32\,M cycles plus lm\_head:
\SI{0.64}{s} at 50\,MHz (Nexus FPGA) and \SI{32}{ms} at the 1\,GHz datasheet
clock for the body alone, i.e.\ 1.6 and 31 tokens/s; with a full-vocabulary
lm\_head (a further $\approx$60\,M cycles) 0.5 and 11 tokens/s. These are the numbers the GR-288 design point
(16\,384 MAC/cycle, model-resident SRAM) is meant to improve by three orders of
magnitude; the per-op table above is the baseline it is measured against.

\section{Simulator throughput}
\subsection{Diagnosis}
The first runs on a c7i.2xlarge (8 vCPU) EC2 host executed T1 in 300\,s of
wall time for 2.13\,M cycles: 7.2\,k cycles/s, and hence about 90 minutes per
token. Verilator itself is not that slow for this design. The cost is in the
cocotb testbench: (i) the AXI bus monitor and the halt poller each wake a
Python coroutine on every clock edge and read half a dozen VPI signals;
(ii) every 16-byte weight beat requested by the core is served by a chain of
Python coroutines (address monitor, memory model, data driver), and a token
pulls 100\,MB of weights, i.e.\ $6.25$\,M beats. Removing the idle-cycle
polling alone would leave the beat path, which at $\sim$100\,\si{\micro\second}
per beat is still tens of minutes per token.

\subsection{Fix: external memory in the RTL}
We added \texttt{axi\_sim\_mem.sv}, a simulation-only AXI4 slave whose storage
lives in C++ (\texttt{ddr\_sim\_mem.cc}) behind DPI, and a generated wrapper
top \texttt{RvvCoreMiniHighmemAxiSimMem} that instantiates the emitted core and
connects its AXI master to that memory; all other ports pass through
unchanged. The test loads and inspects the memory through exported C
functions (\texttt{ddr\_backdoor\_*\_c}, via \texttt{ctypes}, the same
mechanism as the existing SRAM backdoor), so weight traffic never enters
Python. In the cocotb interface, \texttt{wait\_for\_halted} now blocks on the
rising edge of \texttt{io\_halted} with a \texttt{Timer} timeout, and the bus
monitor sleeps on an \texttt{Event} whenever no testbench-initiated transaction
is in flight. The wrapper is detected automatically (no
\texttt{io\_axi\_master\_*} ports), so the change is backward compatible with
every existing test. Per-cycle Python work during compute is now zero.

One pitfall cost an iteration: the Verilator configuration template marks
every \texttt{io\_*} variable of the top module public, so a wrapper whose
internal nets were named \texttt{io\_axi\_master\_*} still exposed them to
cocotb, the interface concluded the Python memory model was in use, and two
drivers fought over the read-data channel (one AR, two R beats, the first a
SLVERR). Renaming the internal nets fixed it; the interface now also logs which
memory mode is active.

Table~\ref{tab:speed} gives the measured throughput for T1 (2.27\,M
simulated cycles including reset and loading). The RTL memory alone gives
2.3$\times$; the cocotb clock is already C-implemented, and profiling
attributes the remaining $\sim$60\,\si{\micro\second} per cycle to
Verilator's evaluation of the core plus the hand-written vector backend.
Verilator's thread scheduler rejects \texttt{--threads 4} for this design
(\texttt{UNOPTTHREADS}); forcing threads anyway gives the last two rows.

\begin{table}[t]
\centering\small
\caption{Simulator throughput for T1 on c7i.2xlarge (3.2\,GHz).}
\label{tab:speed}
\begin{tabular}{lrrr}
\toprule
Configuration & ns/s & cycles/s & min/token \\
\midrule
stock model, Python AXI memory & \num{9014} & 7.2\,k & 89 \\
RTL memory, event-driven cocotb & \num{20660} & 16.5\,k & 32 \\
\quad + \texttt{--threads 2} & \num{30960} & 24.8\,k & 22 \\
\quad + \texttt{--threads 4} & \num{32384} & 25.9\,k & 21 \\
\quad + \texttt{-O3 --x-initial fast -march=native} (2 thr., A/B) & \num{33026} vs \num{30217} & +9\,\% & 20 \\
\bottomrule
\end{tabular}
\end{table}

\subsection{Machine choice}
Two Verilator threads are the efficient operating point: 1.5$\times$ for
2 cores, while four threads add 4\,\% for twice the cores. Four test tiers
running concurrently at two threads each therefore fill the existing
c7i.2xlarge (8 vCPU, \$0.36/h) exactly; a c7i.4xlarge (\$0.71/h) would only
help if more than four simulations run at once, and a c7a.2xlarge (3.7\,GHz
versus 3.2\,GHz, \$0.41/h) would buy roughly 15\,\% single-thread speed for a
15\,\% higher rate plus a re-bootstrap of the host. Neither changes the
conclusion that a token is tens of minutes in RTL simulation, so the instance
was kept. The gains that matter came from removing Python from the per-cycle
and per-beat paths (3.4$\times$ end to end at two threads); anything beyond
that needs the FPGA (about \SI{0.65}{s} per token at 50\,MHz) or a faster
simulator kernel. Arcilator (CIRCT) was considered and not pursued:
it has no VPI/cocotb integration, and the vector backend is hand-written
SystemVerilog with DPI hooks whose coverage under \texttt{circt-verilog} is
unverified; the Verilator fix above recovers the dominant loss without
changing simulators.

\section{Limitations and next steps}
Positions stay below 512, so sliding-window layers are exercised as full
attention. Prefill is sequential decode; a causal-tiled prefill kernel would
recover the 2-D matmul paths. Activation quantisation is per token and per
tensor; if T4/T6 show logits drifting from HF, per-group scales on the
down-projection input are the first thing to try, and both references go
through one function each. Three different $\exp$ approximations coexist in the
kernel set; the $2\times10^{-2}$ strict tolerance exists to absorb them and
should be tightened from the measured margins. The lm\_head dominates decode
cost at this MAC rate; a model-resident, higher-throughput GEMV is the largest
single lever, exactly as the GR-288 spec assumes.

\section*{Reproduction}
\begin{lstlisting}
bazel run //tests/cocotb/rvv/ml_ops/gemma_inference:dump_gemma3_model -- \
  --model_id unsloth/gemma-3-270m-it --out_dir .../test_data --full_lm_head
bazel test //tests/cocotb/rvv/ml_ops/gemma_inference:gemma3_inference_cocotb_test_gemma3_layer0_decode_step
bazel test //tests/cocotb/rvv/ml_ops/gemma_inference:gemma3_inference_cocotb_test_gemma3_next_token_candidates \
  --test_timeout=30000 --test_env=GEMMA3_PROMPT_TOKENS=6
\end{lstlisting}
Host: EC2 c7i.2xlarge, Ubuntu 24.04, Verilator 5.050, Bazel 8.6 via
\texttt{tools/dev/bazel-docker.sh}; reference model
\texttt{unsloth/gemma-3-270m-it} (ungated mirror of
\texttt{google/gemma-3-270m-it}), transformers 5.12.1, torch 2.12 CPU.
\end{document}
